\section{Infraestrutura computacional}

\begin{frame}{As camadas de uma infraestrutura computacional}
  \centering
  \begin{tikzpicture}[
    layer/.style={draw,rounded corners,minimum width=8.5cm,minimum height=0.62cm,align=center},
    flow/.style={-{Latex[length=2mm]},thick}
  ]
    \node[layer,fill=blue!8] (users) {Usuários, serviços e aplicações};
    \node[layer,fill=blue!12,below=3mm of users] (runtime) {Bibliotecas, runtimes e sistemas operacionais};
    \node[layer,fill=blue!16,below=3mm of runtime] (infra) {Computação, armazenamento e virtualização};
    \node[layer,fill=blue!20,below=3mm of infra] (network) {Redes, conectividade e identidade};
    \node[layer,fill=blue!24,below=3mm of network] (physical) {Hardware, firmware e instalações};
    \draw[flow] (users) -- (runtime);
    \draw[flow] (runtime) -- (infra);
    \draw[flow] (infra) -- (network);
    \draw[flow] (network) -- (physical);
  \end{tikzpicture}
  \vspace{0.5em}
  \small\placeholder{Modelo analítico inicial; adaptar ao caso e às fontes escolhidas.}
  \note{A separação em camadas ajuda a localizar dependências, mas as relações reais são muitas vezes horizontais e circulares. O desenho não representa uma pilha universal.}
\end{frame}
